\documentclass[11pt]{article}
\usepackage[a4paper,margin=1in]{geometry}
\usepackage{amsmath,amssymb,amsthm,booktabs,enumitem,mathtools,microtype}
\usepackage[T1]{fontenc}
\usepackage{lmodern}
\usepackage[hidelinks]{hyperref}

\newtheorem{theorem}{Theorem}[section]
\newtheorem{proposition}[theorem]{Proposition}
\newtheorem{lemma}[theorem]{Lemma}
\newtheorem{corollary}[theorem]{Corollary}
\theoremstyle{definition}
\newtheorem{definition}[theorem]{Definition}
\theoremstyle{remark}
\newtheorem{remark}[theorem]{Remark}

\newcommand{\mob}{\mu}
\newcommand{\TT}{\{-1,0,1\}}
\newcommand{\dd}{\delta}
\newcommand{\thh}{\vartheta}
\newcommand{\kap}{\kappa_2}
\newcommand{\Z}{\mathbb Z}
\newcommand{\Ecal}{\mathcal E}

\hypersetup{
  pdftitle={Phasewise Chowla-Free Lag-Two Memory: Exact Max-Plus Optimization and the All-Clock Supremum},
  pdfauthor={RH research program},
  pdfsubject={Exact fixed-clock max-plus optimization for phasewise Chowla-free lag-two memory},
  pdfkeywords={Mobius function, squarefree progressions, max-plus optimization, finite clocks}
}

\title{Phasewise Chowla-Free Lag-Two Memory:\\
Exact Max-Plus Optimization and the All-Clock Supremum}
\author{RH research program}
\date{August 7, 2026}

\begin{document}
\maketitle

\begin{abstract}
Fix a finite clock $q$ before taking $N\to\infty$.  We optimize all
universally distance-two-safe phasewise lag-two tables
\[
 \epsilon_n=f_{n\bmod q}\bigl(\mob(n-2),\mob(n)\bigr)
\]
whose interpolation coefficient $c_{11}(r)$ vanishes at every phase.  This
is the phasewise Chowla-free class, not unrestricted memory.  Exhausting all
$512$ local truth tables leaves exactly $192$ with $c_{11}=0$.  A
subset-dominance reduction first maps them to four relations $0,J,K,I$;
only afterward may $K$ be replaced by $I$, because their full compatibility
profiles agree against canonical neighbors.  The resulting exact
fixed-clock optimum $G(q)$ is a three-state cyclic max-plus problem with
arithmetic weights in $\mathbb Q\pi^{-2}+\mathbb Q\kap$.

At the square clock $q=36$ this gives the exact strict same-clock gain
\[
 G(36)=\frac{9}{2\pi^2}-\frac{\kap}{7}
 >F(36)=\frac4{\pi^2}.
\]
More generally, for the RH-374 clocks $q_y=4\prod_{i\le y}p_i^2$, an exact
run recurrence yields
\[
 G(q_y)=B_y+\Delta_y,
 \qquad
 \Delta_y=\Ecal_y\left(\frac4{A_y\pi^2}-\frac{\kap}{D_y}\right)>0,
 \qquad \Delta_y\to0,
\]
where $\Ecal_y$ counts even-length positive runs.  A separate arbitrary-clock
argument lifts fixed $q$ to $Q_y=\operatorname{lcm}(q,q_y)$, retains a
one-site independent set, and charges every discarded $J$ contribution to
$\sum_{p>p_y}p^{-2}$.  It proves
\[
                    \sup_{q<\infty}G(q)=B_\infty.
\]
No finite-clock attainment or nonattainment statement, growing clock,
phasewise $c_{11}\ne0$ theorem, adaptive-capacity limit, operator, trace,
zero model, or implication for the Riemann hypothesis is asserted.
\end{abstract}

\noindent\textbf{Keywords:} M\"obius function; squarefree progressions;
max-plus optimization; cyclic independent sets; fixed finite clocks.

\section{Problem, class, and exact arithmetic}

The frozen distance-two capacity allows sign words with no two $+1$ values
at distance two.  RH-375 identifies the all-clock supremum $B_\infty$ only
for universally safe one-site factors \cite{RH375}.  RH-378 classifies the
unphased lag-two tables and isolates the coefficient of the ordinary
shift-two correlation \cite{RH378}.  Here the clock is phase dependent, but
that coefficient is required to vanish separately at each phase.

\begin{definition}[Phasewise Chowla-free lag-two class]
Fix $q\ge1$.  For each $r\in\Z/q\Z$, let
\[
             f_r:\TT^2\longrightarrow\{-1,+1\}.
\]
On every ternary input $a=(a_n)$, with zero padding before its first site,
the output is $\epsilon_n=f_{n\bmod q}(a_{n-2},a_n)$.  The family is
\emph{universally safe} if $\epsilon_n=\epsilon_{n+2}=+1$ never occurs for
any ternary input.  It is \emph{phasewise Chowla-free} if the coefficient
$c_{11}(r)$ in
\begin{equation}
 zf_r(x,z)=c_{01}(r)z+c_{02}(r)z^2+c_{11}(r)xz+c_{12}(r)xz^2
            +c_{21}(r)x^2z+c_{22}(r)x^2z^2
 \label{eq:interpolation}
\end{equation}
is zero for every $r$.
\end{definition}

The displayed monomials form a current-zero interpolation basis on $\TT^2$,
and the expansion in that basis is unique, as proved in RH-378.  The
condition is phasewise: cancellation among
unknown correlations carrying different phase weights is not assumed.

Let
\begin{align}
 \dd_{q,r}&=\lim_{N\to\infty}\frac1N
   \sum_{\substack{n\le N\\n\equiv r\ (q)}}\mob(n)^2,
 \label{eq:delta-def}\\
 \thh_{q,r}&=\lim_{N\to\infty}\frac1N
   \sum_{\substack{n\le N\\n\equiv r\ (q)}}
   \mob(n-2)^2\mob(n)^2.
 \label{eq:theta-def}
\end{align}
The first two padded sites change neither limit.  The progression sieve
gives exact local formulae.  If $p^a\Vert q$, then
\begin{equation}
 \pi^2\dd_{q,r}=
 \begin{cases}
 0,&a\ge2\text{ and }r\equiv0\pmod{p^2}\text{ for some }p,\\[1mm]
 \displaystyle\frac6q\prod_{p\mid q}(1-p^{-2})^{-1}
  \prod_{\substack{p\parallel q\\p\mid r}}(1-p^{-1}),&\text{otherwise}.
 \end{cases}
 \label{eq:delta-local}
\end{equation}
Put $\kap=\prod_p(1-2/p^2)$.  Then
\begin{equation}
 \thh_{q,r}=\beta_{q,r}\kap,
 \label{eq:theta-local}
\end{equation}
where $\beta_{q,r}=0$ if $a\ge2$ and
$r\equiv0$ or $2\pmod{p^2}$ for some $p$, or if $2\parallel q$ and $r$ is
even.  In every other case,
\begin{equation}
 \beta_{q,r}=\frac1q\prod_{p\mid q}(1-2p^{-2})^{-1}
 \prod_{\substack{p\parallel q,\ p\text{ odd}\\r\equiv0\text{ or }2\ (p)}}
 (1-p^{-1}).
 \label{eq:beta-local}
\end{equation}
Both formulae follow by conditioning the squarefree and two-squarefree
Euler products on one residue class \cite{Mirsky1948}.  In particular,
\begin{equation}
                       0\le\thh_{q,r}\le\dd_{q,r}.
 \label{eq:density-cone}
\end{equation}

\subsection{Fixed-clock cancellation}

\begin{proposition}[Phasewise limiting score]
\label{prop:phase-limit}
For every fixed finite $q$ and every phasewise Chowla-free family,
\begin{equation}
 \lim_{N\to\infty}\frac1N\sum_{n\le N}
 \mob(n)f_{n\bmod q}(\mob(n-2),\mob(n))
 =\sum_{r\bmod q}\bigl(c_{02}(r)\dd_{q,r}
                         +c_{22}(r)\thh_{q,r}\bigr).
 \label{eq:phase-limit}
\end{equation}
\end{proposition}

\begin{proof}
Insert \eqref{eq:interpolation}.  The $c_{02}$ and $c_{22}$ terms converge
by \eqref{eq:delta-def}--\eqref{eq:theta-def}, and the $c_{11}$ term is
absent at every phase.  The $c_{01}$ term is a M\"obius sum in one fixed
arithmetic progression and is $o(N)$ by Davenport cancellation and finite
Fourier inversion \cite{Davenport1937}.

For a masked term, use
\[
                       \mob(m)^2=\sum_{d^2\mid m}\mob(d).
\]
Fix a cutoff $R$.  Combining the phase congruence modulo $q$ with
$d^2\mid n$ or $d^2\mid n-2$ produces only finitely many progressions modulo
$\operatorname{lcm}(q,d^2)$ for $d\le R$; Davenport applies to each fixed
progression.  The absolute tail is
\[
 \sum_{R<d\le\sqrt{N+2}}\left(\frac{N}{d^2}+1\right)
       =O(N/R+\sqrt N).
\]
Thus $q$ and $R$ are fixed first, then $N\to\infty$, and only afterward
$R\to\infty$.  Both $c_{12}$ and $c_{21}$ contributions are $o(N)$.
\end{proof}

The first missing extension is now explicit.  If any $c_{11}(r)\ne0$, the
proof meets a phase-weighted shift-two sum
\[
 \sum_{n\le N}c_{11}(n\bmod q)\mob(n-2)\mob(n),
\]
for which the frozen sources supply no natural-average cancellation.  RH-376
identifies even the unweighted $q=1$ case as the ordinary shift-two
Chowla boundary \cite{RH376}.

\section{The 512-table census and canonical dominance}

For a local table $f$, write
\[
              E_f=\{(x,z)\in\TT^2:f(x,z)=+1\}.
\]
Two phase tables $E$ at $r$ and $E'$ at $r+2$ are compatible precisely when
there are no $x,z,w\in\TT$ with $(x,z)\in E$ and $(z,w)\in E'$.  Universal
safety is exactly this test at every consecutive pair along addition by two.

The exact interpolation census of all $2^9=512$ truth tables gives
\begin{center}
\begin{tabular}{c|rrrrr}
\toprule
$c_{11}$&$-1$&$-1/2$&$0$&$1/2$&$1$\\
\midrule
count&32&128&192&128&32\\
\bottomrule
\end{tabular}
\end{center}
Thus exactly $192$ tables enter the declared class.  Their nine possible
main pairs, multiplicities, and canonical targets are as follows.
\begin{center}
\begin{tabular}{rrrcl}
\toprule
$c_{02}$&$c_{22}$&count&target&target minus original payoff\\
\midrule
$-1$&$0$&8&$0$&$\dd$\\
$-1$&$1$&32&$0$&$\dd-\thh$\\
$-1$&$2$&8&$K$&$\dd-\thh$\\
$0$&$-1$&16&$0$&$\thh$\\
$0$&$0$&64&$0$&$0$\\
$0$&$1$&16&$K$&$0$\\
$1$&$-2$&8&$J$&$\thh$\\
$1$&$-1$&32&$J$&$0$\\
$1$&$0$&8&$I$&$0$\\
\bottomrule
\end{tabular}
\end{center}
Here the four canonical plus-edge relations and weights are
\begin{align}
 E_0&=\varnothing,&w_r(0)&=0,\nonumber\\
 E_J&=\{(0,+1)\},&w_r(J)&=J_r:=\dd_{q,r}-\thh_{q,r},\nonumber\\
 E_K&=\{(-1,+1),(+1,+1)\},&w_r(K)&=K_r:=\thh_{q,r},\nonumber\\
 E_I&=\TT\times\{+1\},&w_r(I)&=I_r:=\dd_{q,r}.
 \label{eq:canonical-relations}
\end{align}

\begin{proposition}[Canonical subset dominance]
\label{prop:subset-dominance}
For every one of the $192$ tables, the target relation in the preceding
table is a subset of its plus-edge relation, and its limiting payoff is at
least the original payoff for every $0\le\thh\le\dd$.
\end{proposition}

\begin{proof}
The subset assertion is a finite nine-edge check.  The last column gives the
payoff difference.  Every entry is nonnegative by
\eqref{eq:density-cone}.  Replacing a plus-edge set by a subset can only
remove possible composable pairs, so the replacement preserves compatibility
with both neighbors.
\end{proof}

The compatibility matrix after this reduction is
\begin{equation}
\begin{array}{c|cccc}
 &0&J&K&I\\ \hline
0&1&1&1&1\\
J&1&1&0&0\\
K&1&1&0&0\\
I&1&1&0&0
\end{array}.
\label{eq:compatibility-matrix}
\end{equation}
In particular, $K$ and $I$ now have identical full incoming and outgoing
profiles.  Since
\[
                          I_r=J_r+K_r\ge K_r,
\]
every $K$ may be replaced by $I$ without changing compatibility.  The order
of these two operations is essential: $E_K\subset E_I$, so the replacement
need not preserve compatibility against an arbitrary noncanonical neighbor.
The subset reduction must come first.

\section{Exact three-state cyclic max-plus optimization}

After the preceding reductions the action set is
\[
                         \mathcal A=\{0,J,I\}.
\]
Let $C(a,b)$ be the restriction of \eqref{eq:compatibility-matrix} to this
alphabet.  Addition by two partitions $\Z/q\Z$ into one cycle when $q$ is
odd and two cycles when $q$ is even; for $q=1,2$ the components retain their
self-loops.

\begin{definition}[Fixed-clock max-plus value]
For an addition-by-two cycle $(r_0,\ldots,r_{m-1})$, put
\begin{equation}
 \Gamma(r_0,\ldots,r_{m-1})=
 \max_{\substack{a_0,\ldots,a_{m-1}\in\mathcal A\\
 C(a_i,a_{i+1})=1\ \mathrm{cyclically}}}
 \sum_{i=0}^{m-1}w_{r_i}(a_i).
 \label{eq:cycle-value}
\end{equation}
Define $G^+(q)$ as the sum of \eqref{eq:cycle-value} over all components.
\end{definition}

This is an exact finite optimization, not a numerical scan.  For a
conditioned start $a_0=s$, its dynamic program is
\begin{equation}
 V_0^{(s)}(a)=
 \begin{cases}w_{r_0}(s),&a=s,\\-\infty,&a\ne s,\end{cases}
 \qquad
 V_i^{(s)}(b)=w_{r_i}(b)+\max_{a:C(a,b)=1}V_{i-1}^{(s)}(a),
 \label{eq:max-plus-dp}
\end{equation}
followed by $\max_{b:C(b,s)=1}V_{m-1}^{(s)}(b)$ and a maximum over the three
starts.  All weights lie in
$\mathbb Q\pi^{-2}+\mathbb Q\kap$.  The artifact stores them as exact pairs
and uses a certified rational enclosure of $\pi^2\kap$ for comparisons; an
unresolved comparison fails closed.

There is also a useful independent form.  Start from the compatible all-$J$
configuration.  Replacing $J_r$ by $I_r$ forces the predecessor $r-2$ to
change from $J$ to $0$, for net gain
\begin{equation}
                 I_r-J_r-J_{r-2}=K_r-J_{r-2}.
 \label{eq:vertex-gain}
\end{equation}
Two chosen $I$ sites cannot be adjacent on the addition-by-two cycle.
Conversely, any zero action not immediately preceding an $I$ may be changed
to $J$: its weight is nonnegative and the $J$ column is compatible with
every possible predecessor left by the construction.  Thus an optimum is
encoded exactly by its independent set of $I$ sites.
Consequently
\begin{equation}
 G^+(q)=\sum_{r\bmod q}J_r+
 \operatorname{MWIS}_{+2}\bigl(K_r-J_{r-2}\bigr),
 \label{eq:independent-reduction}
\end{equation}
with self-loops retained.  The executable certificate checks
\eqref{eq:max-plus-dp} and \eqref{eq:independent-reduction} independently.

\begin{lemma}[Input reflection]
If $E_r$ is a universally safe phase family, then
\[
                E_r^*=\{(-x,-z):(x,z)\in E_r\}
\]
is universally safe.  For $c_{11}(r)=0$, its main coefficients
$(c_{02}(r),c_{22}(r))$ are the negatives of the original pair.
\end{lemma}

\begin{proof}
A composable pair in the reflected family negates to a composable pair in
the original family, and conversely.  Also
$zf_r^*(x,z)=-[zf_r(x,z)]_{(x,z)\mapsto(-x,-z)}$.  The even-degree main
coefficients therefore change sign, while $c_{11}=0$ remains zero.  The
terms outside the main pair vanish by Proposition~\ref{prop:phase-limit}.
\end{proof}

Thus the negative of every limiting score is also attainable.  We henceforth
define
\begin{equation}
 \boxed{\displaystyle
 G(q)=\max_{(f_r)}\left|\lim_{N\to\infty}\frac1N
 \sum_{n\le N}\mob(n)f_{n\bmod q}(\mob(n-2),\mob(n))\right|
 =G^+(q).}
 \label{eq:G-definition}
\end{equation}

For orientation, the smallest clocks already show that memory changes the
one-site problem:
\begin{equation}
                 G(1)=G(2)=\frac6{\pi^2}-\kap>0,
 \label{eq:q1}
\end{equation}
whereas $F(1)=F(2)=0$ because their one-site phase graphs have self-loops.

\section{Exact square-clock refinement}

Let $3=p_1<p_2<\cdots$ be the odd primes and put
\begin{equation}
 P_y=\prod_{i\le y}p_i^2,
 \quad q_y=4P_y,
 \quad A_y=\prod_{i\le y}(p_i^2-1),
 \quad D_y=\prod_{i\le y}(p_i^2-2).
 \label{eq:square-parameters}
\end{equation}
On one period of the RH-374 odd squarefree-support word, let $R_j^{(y)}$ be
the number of positive runs of exact length $j$.  All runs have length at
most eight, and RH-374 gives their exact Euler-product formula
\cite{RH374}.  Retain its odd-run count $O_y$ and define the new symbol
\begin{equation}
 O_y=R_1^{(y)}+R_3^{(y)}+R_5^{(y)}+R_7^{(y)},
 \qquad
 \Ecal_y=R_2^{(y)}+R_4^{(y)}+R_6^{(y)}+R_8^{(y)}.
 \label{eq:run-counts}
\end{equation}
The notation $\Ecal_y$ counts even-length runs; it is distinct from the
Euler products and from the number of sites inside even runs used in
RH-374.
The locked one-site square-clock value is
\begin{equation}
                 B_y=\frac{4+2O_y/A_y}{\pi^2},
                 \qquad B_y\longrightarrow B_\infty,
 \label{eq:B-y}
\end{equation}
by RH-374 and RH-375 \cite{RH374,RH375}.

Every positive current phase has
\begin{equation}
             a_y=\frac{2}{A_y\pi^2},
 \label{eq:a-weight}
\end{equation}
and every phase at which both current and predecessor supports are positive
has
\begin{equation}
             b_y=\frac{\kap}{2D_y}.
 \label{eq:b-weight}
\end{equation}
Define $H_y=\kap A_y/D_y$.  The Euler products give
\begin{equation}
 H_y=\frac12
 \prod_{i\le y}\left(1-\frac1{p_i^2}\right)
 \prod_{\substack{p>p_y\\p\text{ odd}}}
        \left(1-\frac2{p^2}\right)
 \nearrow \frac12\prod_{p\text{ odd}}(1-p^{-2})=\frac4{\pi^2}.
 \label{eq:H-limit}
\end{equation}
Moreover $H_y\ge H_1=8\kap/7>2/\pi^2$.  One elementary lower bound is
\[
 \kap=\frac12\prod_{p\text{ odd}}(1-2p^{-2})
 >\frac12\left(1-2\sum_{\substack{n\ge3\\n\text{ odd}}}n^{-2}\right)
 =\frac32-\frac{\pi^2}{8}>\frac14;
\]
here $3<\pi<\sqrt{10}$ suffices.  Together with
\eqref{eq:H-limit}, this proves
\begin{equation}
                         \frac{a_y}{2}<b_y<a_y.
 \label{eq:a-b-order}
\end{equation}

Consider a positive odd-support run of length $n$, bracketed on the left by
a zero phase.  If $M_n$ is its optimal three-state score, then
\begin{equation}
 M_0=0,\quad M_1=a_y,\qquad
 M_n=\max\{M_{n-2}+a_y,\ M_{n-1}+a_y-b_y\}.
 \label{eq:run-recurrence}
\end{equation}
The first alternative ends with $I$ and forces the preceding site to zero;
the second ends with $J$.  Induction using \eqref{eq:a-b-order} gives
\begin{equation}
 M_{2m}=(m+1)a_y-b_y\quad(m\ge1),
 \qquad
 M_{2m+1}=(m+1)a_y\quad(m\ge0).
 \label{eq:run-solution}
\end{equation}
Thus an odd-length run matches its one-site MWIS score, while each
even-length run gains exactly $a_y-b_y$.  The odd phase cycle contains two
copies of the support word; the even phase cycle adds no such correction.

\begin{theorem}[Exact square-clock memory formula]
\label{thm:square-memory}
For every $y\ge1$,
\begin{equation}
 \boxed{\displaystyle
 G(q_y)=B_y+\Delta_y,\qquad
 \Delta_y=\Ecal_y\left(\frac4{A_y\pi^2}-\frac{\kap}{D_y}\right)>0.}
 \label{eq:square-memory}
\end{equation}
Furthermore $\Delta_y\to0$, and hence $G(q_y)\to B_\infty$.
\end{theorem}

\begin{proof}
Equation \eqref{eq:run-solution}, summed over the two copies of every run,
adds $2\Ecal_y(a_y-b_y)$ to the RH-374 one-site value $B_y$, which is
exactly \eqref{eq:square-memory}.  RH-374 proves that an exact length-eight
run persists, so $\Ecal_y>0$; inequality \eqref{eq:a-b-order} proves
strictness.

Every even run contains at least two positive sites and distinct runs are
disjoint, so $\Ecal_y/A_y\le1/2$.  By \eqref{eq:H-limit},
\[
 \Delta_y=\frac{\Ecal_y}{A_y}
           \left(\frac4{\pi^2}-H_y\right)\longrightarrow0.
\]
Finally $B_y\to B_\infty$ by RH-374.  No monotonicity of $\Delta_y$ is used
or asserted.
\end{proof}

For $y=1$, one has $q_1=36$, $A_1=8$, $D_1=7$, $O_1=0$, and
$\Ecal_1=1$.  Therefore
\begin{equation}
 G(36)=\frac9{2\pi^2}-\frac{\kap}{7},
 \qquad
 G(36)-F(36)=\frac1{2\pi^2}-\frac{\kap}{7}.
 \label{eq:q36-value}
\end{equation}
The sign is exact:
\begin{equation}
 \kap<\prod_{p\in\{2,3,5,7\}}(1-2p^{-2})
 =\frac{7567}{22050}<\frac7{20}<\frac7{2\pi^2},
 \label{eq:q36-strict}
\end{equation}
where the last step uses $\pi^2<10$.  Thus \eqref{eq:q36-value} is an
exact square-clock strict gain.  It is not called the first same-clock gain,
because \eqref{eq:q1} already gives $G(1)>F(1)$.

\section{The all-clock supremum}

Let $F(q)$ be the RH-375 one-site optimum.  Every one-site family embeds in
the present class by
\begin{equation}
                         f_r(x,z)=g_r(z).
 \label{eq:one-site-embedding}
\end{equation}
Its score expansion has $c_{11}(r)=0$, and its universal safety is unchanged.
Hence
\begin{equation}
                         G(q)\ge F(q)\quad\text{for every }q.
 \label{eq:lower-embedding}
\end{equation}

The opposite supremum inequality needs a different argument from
Theorem~\ref{thm:square-memory}; no same-prime-support saturation law is
claimed for memory tables.

\begin{theorem}[Exact all-clock phasewise-memory supremum]
\label{thm:all-clock}
Over all fixed finite clocks and universally safe phasewise Chowla-free
lag-two tables,
\begin{equation}
                  \boxed{\displaystyle\sup_{q<\infty}G(q)=B_\infty.}
 \label{eq:all-clock}
\end{equation}
\end{theorem}

\begin{proof}
Fix an arbitrary finite $q$ and an optimal canonical action family for
$G(q)$.  Choose $y$ so that $p_1,\ldots,p_y$ contain every odd prime divisor
of $q$, and put
\begin{equation}
                  Q_y=\operatorname{lcm}(q,q_y).
 \label{eq:cofinal-clock}
\end{equation}
Lift the $q$-periodic family to $Q_y$.  The lifted family produces the same
output word and the same fixed-clock limit.  The clock $Q_y$ is divisible by
$q_y$ and has exactly the same prime support as $q_y$.

On the fine $Q_y$ phases, retain every positive-weight $I$ phase.  Retain a
$J$ phase only when its predecessor phase is forced divisible by $p^2$ for
some supported prime $p\le p_y$, including $p=2$.  At such a $J$ phase,
$\thh=0$, so its contribution $J=\dd$ is a one-site weight.  Delete every
zero-density current phase.

The retained positive phases form a distance-two independent set.  Indeed,
the compatibility column of $I$ in \eqref{eq:compatibility-matrix} forces
its predecessor action to be $0$.  For a retained $J$, the predecessor has
zero squarefree density by construction.  Therefore the total retained
weight is at most the one-site optimum at $Q_y$.  The special RH-375
same-support theorem applies only at this one-site step and gives
\begin{equation}
                   \text{retained contribution}\le F(Q_y)=B_y.
 \label{eq:retained-bound}
\end{equation}

It remains to bound the discarded $J$ phases.  Their payoff
$J=\dd-\thh$ is the density of current-squarefree inputs for which the
predecessor is not squarefree.  Since the predecessor fine phase was not
forced divisible by any supported prime square, every such input has
$p^2\mid n-2$ for some $p>p_y$.  Summing first over the disjoint fine phases,
the number of discarded inputs up to $N$ is at most
\[
 N\sum_{p>p_y}p^{-2}+\pi(\sqrt{N+2})+O(1),
\]
because only primes $p\le\sqrt{N+2}$ can occur and each congruence
$p^2\mid n-2$ has at most $N/p^2+1$ solutions; the $O(1)$ term absorbs the
two zero-padded initial sites.  Division by $N$ and then
$N\to\infty$, with $q,y,Q_y$ fixed, gives
\begin{equation}
              \text{discarded contribution}
              \le\sum_{p>p_y}\frac1{p^2}.
 \label{eq:discarded-bound}
\end{equation}
The clock $q$, the index $y$, and hence $Q_y$ are fixed while the
$N\to\infty$ limits underlying \eqref{eq:retained-bound} and
\eqref{eq:discarded-bound} are taken.  Combining them yields
\[
                 G(q)\le B_y+\sum_{p>p_y}p^{-2}.
\]
Only now send $y\to\infty$.  Since $B_y\to B_\infty$ and the convergent
prime-square tail tends to zero, $G(q)\le B_\infty$ for every fixed $q$.

Conversely, \eqref{eq:lower-embedding} and the RH-375 theorem give
\[
     \sup_qG(q)\ge\sup_qF(q)=B_\infty.
\]
Together the two inequalities prove \eqref{eq:all-clock}.
\end{proof}

The theorem does not decide whether some finite clock attains
$B_\infty$.  The proof supplies neither attainment nor nonattainment for
$G$, and it never chooses $q=q(N)$.

\section{Exact executable certificate}

The standard-library artifact exhausts all local tables.  It evaluates
\eqref{eq:delta-local} and \eqref{eq:beta-local} with exact rational
arithmetic and solves \eqref{eq:max-plus-dp} in the
Euler-symbol basis.  Its deterministic finite fixtures are
\begin{center}
\begin{tabular}{rll}
\toprule
$q$&exact $G(q)$&role\\
\midrule
$1$&$6/\pi^2-\kap$&self-loop boundary\\
$2$&$6/\pi^2-\kap$&two self-loop components\\
$3$&$9/(2\pi^2)-3\kap/7$&small-clock fixture\\
$4$&$4/\pi^2$&small-clock fixture\\
$5$&$15/(4\pi^2)-5\kap/23$&small-clock fixture\\
$6$&$5/\pi^2-3\kap/7$&small-clock fixture\\
$36$&$9/(2\pi^2)-\kap/7$&exact square-clock strict gain\\
$180$&$73/(16\pi^2)-25\kap/161$&composite fixture\\
$900$&$73/(16\pi^2)-\kap/7$&square-clock fixture\\
$44100$&$1177/(256\pi^2)-1105\kap/7567$&square-clock fixture\\
\bottomrule
\end{tabular}
\end{center}
For $q=1,\ldots,6$ and all four larger clocks, the three-state cyclic DP
equals the independent-set reduction \eqref{eq:independent-reduction}.
The square-clock run counts give
\[
        \Ecal_1=1,\qquad \Ecal_2=23,\qquad \Ecal_3=1105,
\]
and reproduce the corresponding $q=36,900,44100$ formulas.  The artifact
also checks $\sum_r\dd_{q,r}=6/\pi^2$ and
$\sum_r\thh_{q,r}=\kap$ for twelve distinct local normalization patterns.

Finite decimal intervals in the result ledger are labelled
\texttt{diagnostic\_only}.  They reproduce the exact rows but play no role
in the cofinal proof or in any asymptotic inference.  The full 512-row
census serialization has SHA-256
\begin{center}
\texttt{aaae39b0af85b13e7cc75baa7170a29f1ac60355443d7b80f3fee06d4af56121}.
\end{center}

\section{Limitations, route verdict, and Gates}

Route A is \texttt{GO}.  The fixed-clock phase cancellation, exhaustive
canonical reduction, exact cyclic max-plus optimizer, square-clock strict
gain, and cofinal union-bound theorem form a standalone exact chain.  Route B
is \texttt{STOP\_SCOPED}.  The first blocker is the phase-weighted
shift-two term displayed after Proposition~\ref{prop:phase-limit}; neither
RH-376 nor any other locked source cancels it in ordinary averages.

The result is limited in four further ways.  First, the exact square-clock
formula is not extended to arbitrary same-support multiples.  Second, no
monotonicity of $\Delta_y$ is claimed.  Third, the finite-clock attainment
question is left open.  Fourth, the scalar supremum is not an ordinary limit
of the adaptive capacity and uses no infinite or growing-clock selector.

Gates A--E remain false/open.  These externally prescribed arithmetic
tables do not define a canonical intrinsic dynamical spectral determinant,
a time-oriented scattering or unitary completion, a self-adjoint generator
with an intrinsic $T\log T$ law, a von-Mangoldt prime-power trace, or a
completed-zeta divisor equality.  Nothing here identifies Riemann zeros,
constructs a Hilbert--P\'olya operator, or proves the Riemann hypothesis.

\section{Conclusion}

The phasewise $c_{11}=0$ memory problem has the same all-finite-clock
supremum as the one-site class, even though memory produces exact strict
gains at individual clocks.  The mechanism is now separated cleanly: local
truth-table dominance gives a three-state optimizer; even square-clock runs
create the finite gain; and a retained one-site set plus an explicit
prime-square tail prevents that gain from changing the cofinal endpoint.
Any larger memory theorem must pay for the phase-weighted $D_2$ terms rather
than infer their cancellation from this finite classification.

\section*{Reproducibility and declarations}

\textbf{Data availability.}  All source code, exact outputs, tests, closed
schema, source hashes, and archive-verification files are contained in the
RH-379 repository directory.  No external dataset was created or analyzed.

\textbf{Ethics.}  This mathematical study uses no human participants,
animals, personal data, or biological materials.  Ethics approval was not
required.

\textbf{Author contributions.}  The RH research program performed
conceptualization, methodology, formal analysis, software, validation, and
writing---original draft and review/editing.

\textbf{Funding.}  This research received no specific grant from any
funding agency in the public, commercial, or not-for-profit sectors.

\textbf{Declaration of interests.}  The authors declare no known competing
financial interests or personal relationships that could have influenced
the work reported here.

\textbf{AI-assisted workflow.}  The manuscript and executable checks were
prepared with AI-assisted academic writing and coding tools under the
repository's documented multi-agent workflow.  The mathematical claims are
bounded by the frozen sources, exact proofs, independent audits, and
reproducible artifacts.  The authors retain responsibility for the final
content and its integrity.

\begingroup
\small
\bibliographystyle{plain}
\bibliography{references}
\endgroup
\end{document}
